\documentclass[11pt,a4paper]{article}
% Shared preamble for RuPhO English problem statements (match T1 2026 style).
% Use: \input{latex/rupho_en_preamble.tex} after \documentclass.
\usepackage[margin=1in]{geometry}
\usepackage{amsmath,amssymb}
\usepackage{graphicx}
\usepackage{float}
\usepackage{enumitem}
\usepackage{microtype}
\usepackage{caption}
\usepackage{tabularx}
\usepackage{hyperref}

\captionsetup{font=small,labelfont=bf,skip=6pt}

\setlist[itemize]{leftmargin=1.2cm}
\setlist[enumerate]{leftmargin=1.2cm}

% One outer box per subproblem: [id | statement | points] (no inner rules)
\newcommand{\problembox}[3]{%
  \par\addvspace{0.42em}%
  \noindent
  \setlength{\fboxsep}{7.5pt}%
  \setlength{\fboxrule}{0.95pt}%
  \fbox{%
    \begin{minipage}{\dimexpr\linewidth-2\fboxsep-2\fboxrule\relax}
    \begin{tabularx}{\linewidth}{@{}>{\raggedright\arraybackslash\bfseries}p{2.1em} @{\hspace{0.9em}} >{\raggedright\arraybackslash}X @{\hspace{0.9em}} >{\raggedleft\arraybackslash\bfseries}p{2.4em}@{}}
    #1 & #2 & #3
    \end{tabularx}
    \end{minipage}%
  }%
  \par\addvspace{0.32em}%
}

\title{\textbf{Parabolic mirror}}
\author{\normalsize X19 (2019) --- English translation}
\date{}
\begin{document}
\maketitle

A vertical cylindrical vessel with radius $R$ is filled with mercury and rotates around the cylinder axis with a constant angular velocity $\omega$. The surface of the liquid is a paraboloid of revolution, and the cross section is given by the relation
$$z=\frac{\omega^2 r^2}{2g},$$
where $z$ is the coordinate along the cylinder axis (the lowest point of the mercury surface is taken as zero), $r$ is the distance to the $z$ axis, and $g$ is the acceleration due to gravity.

An observer looks at the surface of mercury from some point on the $z$ axis. He sees the image of his eye. The cylinder is stopped and the surface of the mercury becomes flat. The observer, without changing his position, again sees the image of his eye. The image turned out to be the same size as before, and it did not flip over.

\problembox{A1}{Find the distance $h$ between the observer's eye and the bottom point of the paraboloid when the vessel was rotating.}{2.00}

\vspace{1em}
\noindent\textit{Source:} \url{https://pho.rs/p/3058}

\end{document}